\documentclass[11pt,a4paper]{article}
\usepackage[utf8]{inputenc}
\usepackage[spanish,es-noquoting]{babel}
\usepackage{amsmath,amssymb,booktabs,geometry,microtype}
\geometry{margin=2.6cm}

\newcommand{\dd}{\mathrm{d}}
\newcommand{\Real}{\mathbb{R}}

\title{Ecuación Proacción, versión 2\\[2pt]
\large Dinámica de segundo orden con resorte magnético fatigable}
\author{}
\date{}

\begin{document}
\maketitle

\section{Estado y actualización}

Un sistema viable tiene $N$ necesidades. El estado de la necesidad $i$ es el par
$(x_i, v_i)$, donde $x_i$ es el \emph{nivel de satisfacción} y $v_i$ su velocidad de cambio.
Cada pulso de duración $\Delta t$ actualiza las dos componentes:

\begin{align}
  v_i(t+\Delta t) &= v_i(t) + a_i(t)\,\Delta t, \label{eq:vel}\\
  x_i(t+\Delta t) &= x_i(t) + v_i(t+\Delta t)\,\Delta t . \label{eq:pos}
\end{align}

La aceleración reúne los cinco términos de la ecuación:

\begin{equation}
  a_i(t) \;=\;
  \underbrace{\lambda_i\!\left(x_i,t\right)}_{\text{deriva basal}}
  \;-\;\underbrace{S_i\!\left(x_i,t\right)}_{\text{resorte magnético}}
  \;-\;\underbrace{c_i\,v_i(t)}_{\text{freno}}
  \;+\;\underbrace{u_i(t)}_{\text{estímulos}}
  \;+\;\underbrace{\sum_{j\neq i} W_{ij}\,\big(x_j(t)-x_j^{*}\big)}_{\text{acoplamiento}} .
  \label{eq:epa}
\end{equation}

Escribimos $d_i(t) = x_i(t) - x_i^{*}$ para el desvío respecto del punto de equilibrio
$x_i^{*}$. Convención de signo: $x_i$ mide \emph{satisfacción}, de modo que $d_i<0$ es
déficit y $d_i>0$ superávit.

La deriva basal se aplica en todo pulso, haya o no estímulos. Es el término que vuelve
al sistema no reactivo: aun con todas las variables observadas dentro de rango, la
presión crece en la dirección basal de cada necesidad.

\section{Los términos}

\subsection{Deriva basal}

\begin{equation}
  \lambda_i(x_i,t) \;=\; \sigma_i \, \lambda_i^{0}\, \varphi_i\!\left(|d_i|\right),
  \qquad \sigma_i \in \{-1,+1\},
\end{equation}

donde $\sigma_i$ es la \emph{dirección basal} ($-1$ decae, $+1$ recupera), $\lambda_i^0>0$
su magnitud, y $\varphi_i$ su \emph{forma}. Los casos usados:

\begin{equation}
  \varphi_i(s) =
  \begin{cases}
    1 & \text{lineal (por defecto)},\\[2pt]
    e^{\,\gamma_i s} & \text{exponencial},\\[2pt]
    1 - e^{-s/s_i} & \text{saturante}.
  \end{cases}
\end{equation}

Con forma lineal la deriva es constante; con forma exponencial se acelera con el desvío
y puede producir un punto de inflexión que la lineal no tiene.

\subsection{Resorte magnético}

El resorte atrae hacia el punto de equilibrio con una fuerza que \emph{pierde agarre con
la distancia} y \emph{se desgasta con el uso}:

\begin{equation}
  S_i(x_i,t) \;=\; \kappa_i^{\mathrm{ef}}(t)\; d_i(t)\; e^{-|d_i(t)|/w_i},
  \qquad
  \kappa_i^{\mathrm{ef}}(t) \;=\; \kappa_i \, e^{-f_i \Lambda_i(t)} .
  \label{eq:spring}
\end{equation}

La \emph{carga alostática} $\Lambda_i$ acumula exposición al desvío, con recuperación
opcional de tasa $\rho_i^{\Lambda}$:

\begin{equation}
  \frac{\dd \Lambda_i}{\dd t} \;=\; |d_i(t)| \;-\; \rho_i^{\Lambda}\,\Lambda_i(t).
\end{equation}

El término $-c_i v_i$ de \eqref{eq:epa} frena la velocidad acumulada. Con $c_i=0$ el
sistema oscila sin amortiguar; valores pequeños de $c_i$ suavizan la oscilación sin
alterar los regímenes de la Sección~\ref{sec:regimenes}.

\subsection{Estímulos y señal de calidad}

Los estímulos entran por dos canales. Los \emph{sostenidos} (trabajo en curso, consumo
continuo) actúan sobre la aceleración:

\begin{equation}
  u_i(t) \;=\; \sum_{s} \alpha_{is}\, \mathbb{1}\!\left[s \text{ activo en } t\right],
  \qquad \alpha_{is} \gtrless 0 .
\end{equation}

Los \emph{impulsivos} (una entrega aprobada, una limpieza de disco) actúan como saltos
instantáneos del nivel:

\begin{equation}
  x_i(t^{+}) \;=\; x_i(t^{-}) + \Delta_{i}^{(k)} .
\end{equation}

Con $\alpha_{is}>0$ o $\Delta_i>0$ el estímulo \emph{sacia}; con signo negativo
\emph{vacía}. Una misma acción tiene un vector de efectos sobre todas las necesidades:
compactar el contexto libera memoria y consume tokens.

La \emph{señal de calidad} $g^{(k)}_i \ge 0$ mide cuánto de la saciedad esperada
efectivamente ocurrió:

\begin{equation}
  g^{(k)}_i \;=\; \frac{\Delta_i^{\text{observado}}}{\Delta_i^{\text{esperado}}} .
\end{equation}

Una acción es \emph{adaptativa} respecto de la necesidad que la originó si $g_i>0$.
La distinción importa: $\Delta^{\text{esperado}}$ es lo que el modelo predice y
$g$ es lo que ocurrió. Una acción puede ejecutarse y no satisfacer nada.

\subsection{Acoplamiento}

$W \in \Real^{N\times N}$, con $W_{ii}=0$. El término $W_{ij}(x_j-x_j^{*})$ traslada el
desvío de una necesidad a la aceleración de otra. Es acoplamiento \emph{por estado}, y se
evalúa en cada pulso; el acoplamiento \emph{por evento} se implementa aparte, mediante
suscripción a los cambios de zona.

\section{Regímenes}
\label{sec:regimenes}

El agarre del resorte, $|S_i|$ como función de $|d_i|$, crece hasta $|d_i|=w_i$ y decae
después. Su valor máximo es

\begin{equation}
  G_i \;=\; \frac{\kappa_i^{\mathrm{ef}}\, w_i}{e} .
\end{equation}

Comparando $G_i$ con la magnitud de la deriva se obtienen dos regímenes cualitativamente
distintos:

\begin{center}
\begin{tabular}{ll}
\toprule
\textbf{Condición} & \textbf{Comportamiento bajo negligencia basal}\\
\midrule
$G_i < \lambda_i^{0}$ & la deriva gana siempre: el sistema cae hasta el límite y muere\\
$G_i > \lambda_i^{0}$ & regulación estable: oscila cerca del equilibrio y no muere solo\\
\bottomrule
\end{tabular}
\end{center}

En el segundo régimen existe un \emph{punto de no retorno} $d_i^{*}$: el desvío a partir
del cual el agarre ya no alcanza para compensar la deriva. Es la raíz lejana de
$\kappa_i^{\mathrm{ef}} d\, e^{-d/w_i} = \lambda_i^{0}$, que en forma cerrada resulta

\begin{equation}
  d_i^{*} \;=\; -\,w_i\, W_{-1}\!\left(-\frac{\lambda_i^{0}}{\kappa_i^{\mathrm{ef}} w_i}\right),
  \label{eq:dstar}
\end{equation}

con $W_{-1}$ la rama inferior de la función $W$ de Lambert. Equivalentemente, dado el
punto de no retorno deseado, el alcance se despeja como

\begin{equation}
  w_i \;=\; \frac{d_i^{*}}{\ln\!\left(\kappa_i^{\mathrm{ef}} d_i^{*} / \lambda_i^{0}\right)} .
  \label{eq:w}
\end{equation}

\paragraph{Regla de diseño.} El punto de no retorno debe caer \emph{dentro} del rango
viable, con margen:
\begin{equation}
  d_i^{*} \;<\; \big|L_i^{-} - x_i^{*}\big| .
\end{equation}
Si $d_i^{*}$ cae sobre el límite de viabilidad o más allá, el alcance finito no produce
ningún efecto y el resorte degenera en uno de agarre siempre creciente. Esta condición
convierte a $w_i$ en un parámetro \emph{derivado} del contrato, no elegido.

\paragraph{Muerte por negligencia.} Con $f_i>0$ el agarre se desgasta: un sistema que
regula al principio pierde capacidad con la exposición sostenida, y la negligencia basal
termina en muerte operativa aunque exista régimen estable. Los dos decaimientos son
complementarios: $w_i$ hace perder agarre \emph{con la distancia}, $f_i$ lo hace perder
\emph{con el tiempo}.

\section{Presión, zonas y viabilidad}

\subsection{Presión}

La presión de la necesidad $i$ combina tres fuentes, que se mantienen trazables por
separado:

\begin{equation}
  p_i \;=\; \max\Big\{\, p_i^{\text{nivel}},\;\; p_i^{\text{ritmo}},\;\; p_i^{\text{aut}} \,\Big\} .
\end{equation}

Las dos primeras provienen de las variables observadas $y_{i1},\dots,y_{ik}$ que la
necesidad vigila, y la tercera de la dinámica autónoma:

\begin{align}
  p_i^{\text{nivel}} &= \max_k \; \frac{|z_{ik}|}{1-|z_{ik}|},
    \qquad z_{ik}=\frac{y_{ik}-y^{*}_{ik}}{L_{ik}-y^{*}_{ik}}, \\[4pt]
  p_i^{\text{ritmo}} &= \max_k \; \left[1 - \frac{\tau_{ik}}{T^{\text{rec}}_{ik}+T^{\text{reac}}_{ik}}\right]_{0}^{1},
    \qquad \tau_{ik}=\frac{\text{margen}_{ik}}{\text{ritmo}_{ik}}, \\[4pt]
  p_i^{\text{aut}} &= \frac{|d_i|}{|L_i-x_i^{*}|} \;+\; \eta_i\,\frac{|v_i|}{v_i^{\text{ref}}} .
\end{align}

El máximo, y no la suma, porque los recursos no son fungibles: holgura en una variable no
compensa el agotamiento de otra. En $p_i^{\text{aut}}$, la velocidad $v_i$ cumple el papel
que la tensión acumulada cumplía en la formulación pulsátil: a igual desvío, un sistema
que viene acelerando hacia el déficit está en peor situación que uno detenido.

\subsection{Zonas algedónicas}

Cuatro zonas —verde, amarillo, ámbar y rojo— definidas sobre $|d_i|$ normalizado por el
límite. El verde está centrado en $x_i^{*}$; amarillo, ámbar y rojo aparecen a cada lado,
en déficit y en superávit.

Cada corte $\theta_z$ es una banda de ancho $2\omega_z$, con pertenencia difusa

\begin{equation}
  \mu_z(x) \in [0,1], \qquad \sum_z \mu_z(x) = 1,
\end{equation}

que permite priorizar entre necesidades que están en la misma zona. Y cada corte tiene
histéresis: se entra en la zona al cruzar $\theta_z^{\text{in}}$ y se sale recién al
cruzar $\theta_z^{\text{out}}$, con $\theta_z^{\text{out}} < \theta_z^{\text{in}}$ para
los cortes de déficit. Sin histéresis, un valor que oscila alrededor de un corte emite un
evento por cada cruce.

\subsection{Conjunto viable}

\begin{equation}
  K \;=\; \Big\{\, x \in \Real^{N} \;:\; L_i^{-} \le x_i \le L_i^{+} \;\; \forall i \,\Big\} .
\end{equation}

La condición es \emph{conjuntiva}: basta que una necesidad quede fuera de su rango para
que el sistema deje de ser viable. Salir de $K$ es la muerte operativa.

\section{Censo de hiperparámetros}

\begin{center}
\begin{tabular}{llp{6.2cm}}
\toprule
\textbf{Macro} & \textbf{Subparámetros} & \textbf{Procedencia}\\
\midrule
Punto de equilibrio & $x^{*}$, $L^{-}$, $L^{+}$ & contrato de viabilidad\\
Deriva basal & $\lambda^{0}$, dirección $\sigma$, forma $\varphi$ & contrato; la dirección por el tipo de necesidad\\
Resorte magnético & $\kappa$, $w$, $f$, $c$ & $w$ derivado por \eqref{eq:w}; $\kappa$ por el régimen; $f$ por el tiempo de muerte por negligencia; $c$ por defecto nulo\\
Estímulos & $\alpha_{is}$, $\Delta_i$, $g_i$ & medibles por acción\\
Acoplamiento & $W$ & medible\\
Zonas algedónicas & $\theta_z$, histéresis, ancho $\omega_z$ & contrato; la histéresis por la granularidad de acción\\
Pulso & $\Delta t$ & operativo\\
\bottomrule
\end{tabular}
\end{center}

Quedan libres, por necesidad, $\lambda^{0}$, $\kappa$ y $f$. Las tres responden a una
pregunta del dominio: cuánto se degrada la necesidad por unidad de tiempo, qué régimen de
regulación se busca, y en cuánto tiempo la negligencia debe resultar fatal.

\section{Categorías de necesidad}

\begin{center}
\begin{tabular}{lll}
\toprule
& \textbf{Empuje} & \textbf{Arrastre}\\
\midrule
Dirección basal & $\sigma=-1$ (decae) & $\sigma=+1$ (recupera)\\
Efecto de actuar & sacia ($\alpha>0$) & vacía ($\alpha<0$)\\
Sin actuar & la presión crece & la presión cede\\
Mueve hacia & la activación & la inhibición\\
\bottomrule
\end{tabular}
\end{center}

Un sistema viable necesita al menos una necesidad de cada categoría. Con solo necesidades
de arrastre, la política óptima es la inacción, porque la inacción siempre las mejora. Con
ambas, no existe un nivel de reposo viable: el equilibrio posible es un \emph{ritmo}.

Inhibir no implica necesariamente que el sistema quede quieto: una necesidad de arrastre
puede activar acciones de preservación (liberar memoria, compactar contexto, degradar el
modelo) que atenúan el consumo de lo que protege. Una acción de inhibición que no deja al
sistema en mejores condiciones de sostener su propósito no es inhibición.

\end{document}
